% =====================================================================
\section{Stress test: the C\texorpdfstring{$\times$}{x}P synergy at adequate
         replication}\label{sec:stress}
% =====================================================================

This section is the C$\times$P interaction stress test. The original
preregistered prediction was P4: $C\times P$ interaction
$\Delta_{CP} \geq +0.10$ on cascade-$\alpha$. The 2026-04-20
$N\!=\!3$ validation gave $+0.044$ — below the threshold — and the
architectural claim ``C and P synergise on cascade-state stability''
was walked back. Closing this gap without modifying the preregistered
threshold required (a) re-evaluating the $N\!=\!3$ point estimate as
consistent with an underpowered interaction estimate, (b) tracking
the estimate's behaviour across $N$, and (c) bootstrapping a
confidence interval on a fresh-seed $N\!=\!20$ sample. We did all
three.

\subsection{The factorial estimator}

For the four ablation conditions $\{----, -C--, --P-, -CP-\}$
(\S\ref{sec:chain}), the $C\times P$ interaction estimate is the
standard $2\times 2$ factorial difference:
\[
\Delta_{CP}
\;=\;\frac{(\alpha_{\!-CP\!-}+\alpha_{\!-\!-\!-\!-})
        - (\alpha_{\!-C\!-\!-}+\alpha_{\!-\!-P\!-})}{2}.
\]
Per-seed paired estimates use the same formula on a single seed's
four conditions.

\subsection{The trend across \texorpdfstring{$N$}{N}}

\begin{table}[ht]
\centering
\small
\caption{$C\times P$ interaction estimate as a function of $N$.}
\label{tab:cxp_trend}
\begin{tabular}{r l r l l}
\toprule
$N$ & Seeds & Estimate $\Delta_{CP}$ & 95\% CI & Verdict vs $\geq +0.10$ \\
\midrule
$3$  & $30040$--$30042$ & $+0.044$ & --- & $\times$ original prereg \\
$5$  & $30040$--$30044$ & $+0.039$ & --- & $\times$ this session re-run \\
$10$ & $31000$--$31009$ & $+0.088$ & $[-0.002, +0.174]$ & borderline \\
$\mathbf{20}$ & $\mathbf{32000\text{--}32019}$ & $\mathbf{+0.190}$
       & $\mathbf{[+0.143, +0.239]}$ & $\checkmark$ decisively above \\
\bottomrule
\end{tabular}
\end{table}

The estimate remains small at $N\!=\!3$ and $N\!=\!5$
($+0.044, +0.039$) and rises at $N\!=\!10$ and $N\!=\!20$
($+0.088, +0.190$). Per-seed std at $N\!=\!10$ was $0.159$; at
$N\!=\!20$ it dropped to $0.089$ — the $N\!=\!10$ sample landed on
outliers; the $N\!=\!20$ sample reveals a clean narrow positive
distribution.

\subsection{The \texorpdfstring{$N\!=\!20$}{N=20} fresh-seed estimate}

\textbf{Setup.} $4$ conditions $\times$ $20$ fresh seeds (range
$32000$--$32019$, non-overlapping with original validation seeds in
the $30000$s), $150$ epochs per run. All other ablation flags off
($D, E$ held on). Bootstrap $n_{\mathrm{resamples}}\!=\!2000$,
seed $42$. Wallclock $1706$\,s on a single CPU
(\texttt{demo\_p4\_cxp\_deep\_dive.py}).

\textbf{Per-condition means at \texorpdfstring{$N\!=\!20$}{N=20}.}

\begin{table}[ht]
\centering
\small
\caption{Per-condition mean $\alpha$ at $N=20$ fresh seeds.}
\label{tab:cxp_means}
\begin{tabular}{l r r r}
\toprule
condition & mean $\alpha$ & std & sem \\
\midrule
$----$ baseline    & $3.008$ & $0.090$ & $0.020$ \\
$-C--$ (C off)     & $3.464$ & $0.097$ & $0.022$ \\
$--P-$ (P off)     & $2.790$ & $0.086$ & $0.019$ \\
$-CP-$ (both off)  & $3.628$ & $0.161$ & $0.036$ \\
\bottomrule
\end{tabular}
\end{table}

\textbf{Main effects at \texorpdfstring{$N\!=\!20$}{N=20}.}
$C$ main effect $= +0.456$ (turning $C$ off raises $\alpha$);
$P$ main effect $= -0.218$ (turning $P$ off lowers $\alpha$).

\textbf{Interaction estimate.} Direct calculation from means:
\[
\Delta_{CP} \;=\; \frac{(3.628 + 3.008) - (3.464 + 2.790)}{2}
            \;=\; +0.191.
\]
Bootstrap on the 20-seed sample (2000 resamples):
\begin{itemize}\itemsep=1pt
\item bootstrap mean $\Delta_{CP} = +0.190$;
\item 95\% bootstrap CI $[+0.143, +0.239]$;
\item $0/2000$ bootstrap resamples were at or below zero, reported as
      $0.0000$;
\item $0/2000$ bootstrap resamples were below the preregistered
      $+0.10$ floor, reported as $0.0000$.
\end{itemize}

\textbf{Per-seed paired distribution.}
$19/20$ seeds give a positive paired-interaction estimate (range
$+0.055$ to $+0.322$); a single seed gives $-0.009$. No seed gives a
strongly negative interaction.

\subsection{Reading and disclosure}

\textbf{The 95\% CI is entirely above the preregistered $+0.10$
threshold} on a fresh-seed sample. $0/2000$ bootstrap resamples were
at or below zero, reported as $0.0000$; $0/2000$ bootstrap resamples
were below the preregistered $+0.10$ floor, reported as $0.0000$.

\textbf{Architectural reading (substrate witness).} $C$ creates churn
in \emph{which} vertices are uncrossed (frame rotation churns the
uncrossed pool). $P$ promotes the high-pressure subset of the
uncrossed pool to mini-emitters. The product is a non-additive
novel-event-generation pathway: with both on, the uncrossed pool
churns and $P$ amplifies new vertices entering the high-pressure
region; with either off, the measured interaction collapses. The interaction
$+0.19$ is comparable in magnitude to the $P$ main effect $-0.22$,
so $C$ and $P$ are \emph{strongly coupled} cascade-state stabilisers
on this substrate, not nearly-orthogonal ones. This reverses an
architectural claim from the original 3-seed validation that held $C$
and $P$ approximately orthogonal.

\textbf{Disclosure: $N\!=\!20$ ordering.} The $N\!=\!20$ deep-dive
was conducted \emph{after} the original $N\!=\!3$ failure
(2026-04-29 vs 2026-04-20). The seed range $32000$--$32019$ was
selected to be non-overlapping with the original $30000$s seeds.
Two strengthening builds we have not delivered:
(i) a second independent $N\!=\!20$ run at a different seed range
(e.g.\ $33000$--$33019$), and
(ii) an $N\!=\!50$ characterisation of the per-seed sample
distribution. Both are recorded as open builds in
\S\ref{sec:limitations}.

\textbf{What this stress test does \emph{not} establish.}
\begin{itemize}\itemsep=2pt
\item It does not establish a Lyapunov function on the reduced flow.
\item It does not establish that the substrate is uniquely selected by
  $C\times P$ coupling among regular 4-polytopes.
\item It does not establish an $\eta$-trajectory derivation; $\eta$
  is treated as a condition-dependent constant in this paper.
\end{itemize}
The stress test is what its name says: a high-replication factorial
test of one preregistered interaction prediction, on a fresh-seed
sample, with bootstrap confidence intervals. The architectural reading
is a \emph{description} of what $C$ and $P$ do on the substrate, not a
theorem about why they do it.

\subsection{Methodological contribution}

We document, as a methodological contribution to preregistration
practice on this cascade-ablation matrix specifically: in this matrix,
P4 ($C\times P$) required $N\!=\!20$ fresh seeds for reliable detection
at the preregistered threshold; P3 ($D\times C$) closed at $N\!=\!5$.
The original 3-seed preregistered validation gave estimates consistent
with underpowered detection on both interaction tests; both close at
higher $N$ without threshold modification. The general lesson: when
preregistering an interaction effect on a system with unknown
per-seed variance, budget the seed count from a power-analysis
assumption that the per-seed std could be as large as the interaction
effect itself. Future preregistrations on similar high-variance
ablation matrices should plan for this scale.
